\documentclass[12pt, a4paper]{article}
\usepackage{../../../myconfig}

\begin{document}

\titlebox{Chủ đề: Oxyz}{Hệ trục toạ độ \\ trong không gian}{Oxyz: Hệ trục toạ độ}

% Câu 1 [NB]
\questionLinesManual{0}{1}{Trong không gian tọa độ $Oxyz$, cho $A(1;1;-2)$ và $B(2;-1;0)$. Tọa độ của vectơ $\vec{AB}$ là:}{%
    \begin{tasks}(4)
        \task $\vec{AB}=(1;-2;2)$
        \task $\vec{AB}=(1;2;2)$
        \task $\vec{AB}=(-1;2;-2)$
        \task $\vec{AB}=(3;0;-2)$
    \end{tasks}
}

% Câu 2 [NB]
\questionLinesManual{0}{2}{Trong không gian $Oxyz$, cho vectơ $\vec{u}=2\vec{i}-5\vec{k}$. Tọa độ của vectơ $\vec{u}$ là:}{%
    \begin{tasks}(4)
        \task $(0;2;-5)$
        \task $(2;0;5)$
        \task $(2;-5;0)$
        \task $(2;0;-5)$
    \end{tasks}
}

% Câu 3 [NB]
\questionLinesManual{0}{3}{Trong không gian với hệ tọa độ $Oxyz$, cho điểm $A(1;2;-1)$. Tọa độ hình chiếu vuông góc của $A$ trên mặt phẳng $(Oyz)$ là:}{%
    \begin{tasks}(4)
        \task $(0;2;-1)$
        \task $(1;0;0)$
        \task $(1;2;0)$
        \task $(1;0;-1)$
    \end{tasks}
}

% Câu 4 [TH]
\questionLinesManual{0}{4}{Trong không gian với hệ tọa độ $Oxyz$, cho hai điểm $A(1;2;3)$, $B(3;8;5)$. Tọa độ trung điểm $I$ của đoạn thẳng $AB$ là:}{%
    \begin{tasks}(4)
        \task $I(2;6;2)$
        \task $I(1;3;1)$
        \task $I(4;10;8)$
        \task $I(2;5;4)$
    \end{tasks}
}

% Câu 5 [TH]
\questionLinesManual{0}{5}{Trong không gian $Oxyz$, cho tam giác $ABC$ với $A(1;3;4)$, $B(2;-1;0)$ và $C(3;1;2)$. Toạ độ trọng tâm $G$ của tam giác $ABC$ là:}{%
    \begin{tasks}(4)
        \task $G\left(3;\dfrac{2}{3};3\right)$
        \task $G(2;-1;2)$
        \task $G(2;1;2)$
        \task $G(6;3;6)$
    \end{tasks}
}

% Câu 6 [TH]
\questionLinesManual{0}{6}{Trong không gian $Oxyz$, cho $\vec{a}=(2;-3;3)$, $\vec{b}=(0;2;-1)$, $\vec{c}=(3;-1;5)$. Tọa độ của véctơ $\vec{u}=2\vec{a}+3\vec{b}-2\vec{c}$ là:}{%
    \begin{tasks}(4)
        \task $(-2;-2;7)$
        \task $(10;-2;13)$
        \task $(-2;2;-7)$
        \task $(-2;2;7)$
    \end{tasks}
}

% Câu 7 [TH]
\questionLinesManual{0}{7}{Trong không gian $Oxyz$, cho điểm $A(-5;2;3)$ và $B$ là điểm đối xứng với $A$ qua trục $Oy$. Độ dài đoạn thẳng $AB$ bằng:}{%
    \begin{tasks}(4)
        \task $\sqrt{34}$
        \task $2\sqrt{38}$
        \task $2\sqrt{34}$
        \task $\sqrt{38}$
    \end{tasks}
}

% Câu 8 [TH]
\questionLinesManual{0}{8}{Trong không gian $Oxyz$, cho hình bình hành $ABCD$ có $A(-3;1;2)$, $B(-2;4;-1)$, $C(1;-3;3)$. Tọa độ điểm $D$ là:}{%
    \begin{tasks}(4)
        \task $(0;-6;6)$
        \task $(-6;7;-2)$
        \task $(2;0;0)$
        \task $(-4;2;5)$
    \end{tasks}
}

% Câu 9 [VD]
\questionLinesManual{0}{9}{Trong không gian $Oxyz$, cho hai véctơ $\vec{a}=(1;-2;1)$, $\vec{b}=(-2;1;1)$. Số đo góc giữa hai véctơ $\vec{a}$ và $\vec{b}$ là:}{%
    \begin{tasks}(4)
        \task $120^\circ$
        \task $60^\circ$
        \task $-30^\circ$
        \task $30^\circ$
    \end{tasks}
}

% Câu 10 [VD]
\questionLinesManual{0}{10}{Trong không gian tọa độ $Oxyz$ cho điểm $M(1;-\sqrt{2};\sqrt{3})$. Tìm điểm $M' \in Ox$ sao cho độ dài đoạn thẳng $MM'$ ngắn nhất.}{%
    \begin{tasks}(4)
        \task $M'(-1;0;0)$
        \task $M'(1;0;0)$
        \task $M'(1;0;\sqrt{3})$
        \task $M'(1;-\sqrt{2};0)$
    \end{tasks}
}

\vfill
\begin{center}
    \textbf{--- HẾT BÀI ĐIỂM DANH ---}
\end{center}

\end{document}